\kuetfrontheading{Abstract}

This thesis asks whether stating a robot manipulator's safety requirement as an explicit
constraint, rather than as a term added to the reward, changes the behaviour of a
reinforcement-learned grasping policy and how predictable that behaviour is across independent
training runs. The question matters because collaborative arms such as the Universal Robots UR5e
now share workspaces with people, and a policy trained by reinforcement learning cannot be
inspected the way a hand-derived controller can; its behaviour is known only from the statistics
of what it does. A shaped reward gives the designer one lever for a safety requirement, a penalty
weight fixed before training and validated only by trial and error. This work instead trains a
UR5e in simulation to pick up and lift a cube, comparing an unconstrained proximal policy
optimisation (PPO) baseline against a constrained variant, PPO-Lagrangian, trained under an
explicit floor on manipulability, across five independent seeds and two constraint budgets, with
an added control arm that isolates the constraint's own effect from the extra network a Lagrangian
method introduces. The control arm reproduced the unconstrained baseline's trained weights
exactly, so every measured difference is attributable to the constraint alone. The constrained
agent held task performance steady at both budgets while cutting true singularity crossings by a
factor of up to forty and eliminating all observed joint-limit contact. Its larger effect, and one
this thesis did not anticipate, was on consistency: the unconstrained baseline's safety cost
varied more than twentyfold across identically trained seeds, a spread the constraint collapsed to
roughly two to five times, at no measurable cost to the grasping task. A manipulator whose safety
requirement is a stated, checkable number rather than a hand-tuned weight is both a more reliable
and a more defensible object to build a real deployment around.

\clearpage
